% !TeX root = ../../../main.tex
\begin{ChTwoAbstract}
We consider the problem of predicting how the likelihood of an outcome of interest for a patient changes over time as we observe more of the patient's data.
To solve this problem, we propose a supervised contrastive learning framework
that learns an embedding representation for each time step of a patient time series.
Our framework learns the embedding space to have the following properties:
(1)~nearby points in the embedding space have similar predicted class probabilities,
(2)~adjacent time steps of the same time series map to nearby points in the embedding space,
and (3)~time steps with very different raw feature vectors map to far apart regions of the embedding space.
To achieve property (3), we employ a \emph{nearest neighbor pairing} mechanism in the raw feature space. 
This mechanism also serves as an alternative to ``data augmentation'', a key ingredient of contrastive learning, which lacks a standard procedure that is adequately realistic for clinical tabular data, to our knowledge.
We demonstrate that our approach outperforms state-of-the-art baselines in predicting mortality of septic patients (MIMIC-III dataset) and tracking progression of cognitive impairment (ADNI dataset). Our method also consistently recovers the correct synthetic dataset embedding structure across experiments, a feat not achieved by baselines. 
Our ablation experiments show the pivotal role of our nearest neighbor pairing.
\end{ChTwoAbstract}
\begin{ChTwoKeywords}
contrastive learning, time series analysis, nearest neighbors
\end{ChTwoKeywords}




 


